% ===========================================================================
\section{Riesgos y su control}
\label{sec:riesgos}

\diagrama{control_riesgos.png}{Control y riesgos}{La unidad de control, la unidad de forwarding
y la unidad de detección de riesgos sobre el pipeline, con los recursos compartidos. Fuente
editable: \repo{docs/diagramas/control_riesgos.drawio}.}{fig:control_riesgos}

\begin{center}
\footnotesize
\begin{tabular}{@{}p{2.6cm} p{4.4cm} p{1.9cm} p{3.6cm} p{1.4cm}@{}}
  \toprule
  \textbf{Riesgo} & \textbf{Ejemplo} & \textbf{Se detecta en} & \textbf{Se resuelve con} & \textbf{Costo} \\
  \midrule
  Datos, distancia 1 o 2 & \texttt{add x1} y luego \texttt{sub \ldots, x1} & EX & forwarding & 0 \\
  Datos, distancia 3     & \texttt{add x1} y tres después \texttt{or \ldots, x1} & ID & bypass del banco & 0 \\
  Carga-uso              & \texttt{lw x1} y luego \texttt{add \ldots, x1} & ID & 1 stall + forwarding & 1 ciclo \\
  Control                & \texttt{beq}/\texttt{bne} tomado, \texttt{jal}, \texttt{jalr} & EX & flush de IF/ID e ID/EX & 2 ciclos \\
  Estructural            & \textemdash & \textemdash & no hay ninguno que obligue a frenar & 0 \\
  \bottomrule
\end{tabular}
\end{center}

Las dependencias WAR y WAW no son riesgos: las instrucciones avanzan en orden, leen sus
registros en ID (o los reciben por forwarding en EX) y escriben solo en WB. Tampoco hay riesgos
por la memoria de datos: un \texttt{sw} seguido de un \texttt{lw} a la misma dirección acceden en
MEM en el orden del programa.

\subsection{Forwarding}
\label{sec:forwarding}

La unidad de forwarding (\repo{rtl/control/forwarding_unit.v}) es combinacional y genera los
selectores de los muxes de los dos operandos de EX:

\begin{lstlisting}[style=verilog,numbers=none]
fwd_a = (ex_mem_reg_write && ex_mem_rd != 0 && ex_mem_rd == id_ex_rs1) ? FWD_EX_MEM :
        (mem_wb_reg_write && mem_wb_rd != 0 && mem_wb_rd == id_ex_rs1) ? FWD_WB     :
                                                                         FWD_ID_EX;
// fwd_b: idem con id_ex_rs2
\end{lstlisting}

\begin{itemize}[nosep]
  \item Si las dos instrucciones anteriores escriben el mismo registro, gana EX/MEM, que tiene el
        valor más reciente.
  \item \texttt{x0} nunca se reenvía, y una burbuja tampoco (tiene \texttt{reg\_write = 0}).
  \item Los operandos ya reenviados alimentan también el comparador del salto y el dato de un
        store, así que un salto que depende de una instrucción aritmética anterior no pierde
        ciclos.
\end{itemize}

\diagrama{riesgo_forwarding.png}{Forwarding}{Las tres distancias de una dependencia de datos: desde
EX/MEM, desde MEM/WB y por el bypass del banco de registros.}{fig:riesgo_fwd}

\subsection{Carga-uso}

El dato de un load sale de la BRAM recién al final de MEM. Si la instrucción siguiente lo
necesita en EX, ese dato todavía no existe y hay que frenarla un ciclo. La detección se hace en ID
mirando la instrucción que está en EX (\repo{rtl/control/hazard_unit.v}):

\begin{lstlisting}[style=verilog,numbers=none]
load_use = id_ex_mem_read && id_ex_rd != 0 && if_id_valid &&
           (id_ex_rd == if_id_rs1 || id_ex_rd == if_id_rs2);
\end{lstlisting}

La comparación es \textbf{conservadora} (\decision{011}): usa los bits de \texttt{rs1} y
\texttt{rs2} de la palabra sin mirar si la instrucción los usa. Puede frenar de más
(\texttt{lui}, \texttt{jal} o el HALT, cuyos bits en esas posiciones son parte del inmediato),
nunca de menos. En los 22 programas del plan hubo un solo falso positivo en 14 stalls, el de
\repo{halt_stall.s}, que existe para provocarlo. Después del stall el load está en WB y el dato
llega a EX por forwarding desde MEM/WB.

\diagrama{riesgo_carga_uso.png}{Carga-uso}{Un ciclo de stall: PC e IF/ID retienen, ID/EX recibe
una burbuja y el dato llega por forwarding desde MEM/WB.}{fig:riesgo_cu}

\subsection{Saltos}

Los saltos se resuelven en EX, sin predicción (\decision{012}): es equivalente a predecir siempre
``no tomado''. Cuando \texttt{redirect = 1}, las dos instrucciones que entraron detrás del salto
(en ID y en IF) siguieron el camino $PC + 4$; en el mismo flanco el PC toma
\texttt{redirect\_pc} y las dos se convierten en burbujas. Ninguna llegó a escribir nada, porque
los efectos ocurren en MEM y WB.

\begin{center}
\small
\begin{tabular}{@{}l c@{}}
  \toprule
  \textbf{Caso} & \textbf{Ciclos perdidos} \\
  \midrule
  \texttt{beq}/\texttt{bne} no tomado & 0 \\
  \texttt{beq}/\texttt{bne} tomado, \texttt{jal}, \texttt{jalr} & 2 \\
  Salto que depende de una instrucción aritmética anterior & $+0$ (forwarding al comparador) \\
  Salto que depende de un load inmediatamente anterior & $+1$ (stall de carga-uso) \\
  \bottomrule
\end{tabular}
\end{center}

Resolver en ID ahorraría un ciclo en cada salto tomado, pero agregaría un stall cuando el salto
depende de la instrucción anterior, que es el caso más común en un lazo. La \decision{012} compara
las opciones con la suma de un arreglo: con los saltos en EX, el lazo de \repo{mixto_lazo.s}
(\texttt{lw}, \texttt{add}, \texttt{addi}, \texttt{addi}, \texttt{bne}) tarda 8 ciclos por vuelta.

\diagrama{riesgo_salto.png}{Salto tomado}{El \texttt{beq} se resuelve en EX en el ciclo 3: el PC
toma el destino y las dos instrucciones del camino equivocado se anulan.}{fig:riesgo_salto}

\subsection{Enables y flush}

La unidad de riesgos combina la carga-uso, la redirección, la parada por HALT y el
\texttt{enable} de la Debug Unit en los \emph{enables} y \emph{flush} de cada registro:

\begin{lstlisting}[style=verilog,numbers=none]
stall        = load_use & ~redirect;
en_pc        = enable & ~stall;
en_if_id     = enable & ~stall;      // tambien es el enable de la memoria de programa
en_id_ex     = enable;
en_ex_mem    = enable;
en_mem_wb    = enable;               // tambien es el enable de la memoria de datos
flush_if_id  = redirect | stop_fetch;
flush_id_ex  = redirect | stall;
\end{lstlisting}

\begin{center}
\small
\begin{tabular}{@{}l l l l l@{}}
  \toprule
  \textbf{Situación} & \textbf{PC} & \textbf{IF/ID} & \textbf{ID/EX} & \textbf{EX/MEM, MEM/WB} \\
  \midrule
  Normal                      & avanza & carga & carga & cargan \\
  Stall por carga-uso         & retiene & retiene & burbuja & cargan \\
  Salto tomado                & destino & burbuja & burbuja & cargan \\
  HALT en vuelo               & retiene & burbuja & carga & cargan \\
  Core detenido (\texttt{enable = 0}) & retiene & retiene & retiene & retienen \\
  \bottomrule
\end{tabular}
\end{center}

Con \texttt{enable = 0} nada con estado cambia, y al volver a 1 las señales combinacionales se
recalculan igual: un stall continúa donde quedó, así que se puede avanzar ciclo a ciclo con
\texttt{STEP} en medio de un stall sin perder instrucciones.

\subsection{Riesgos estructurales}

Ningún recurso compartido obliga a frenar. Las memorias de programa y de datos están separadas
(Harvard); el banco tiene dos puertos de lectura y uno de escritura, y el conflicto entre ID y WB
se resuelve con el bypass interno; \texttt{PC+4} y \texttt{PC+imm} tienen sumadores propios; no
hay etapas multiciclo; y la Debug Unit usa puertos propios (el puerto~B de las BRAM y el tercer
puerto del banco), con el core detenido o en reset.
